% F2 --- Task anatomy card for insurance_task5.
% Rendered from scripts/fig02_task_card.html via headless Chrome:
%   chrome --headless --print-to-pdf=assets/task_card.pdf fig02_task_card.html
% Referenced from sections/03_design.tex.

\begin{figure}[t]
  \centering
  \includegraphics[width=\linewidth]{task_card}
  \caption{Anatomy of a \bench{} task. Example: \texttt{insurance\_task5} (Enterprise Property Insurance Claim), a six-turn adjudication of a \textyen1.2\,M fire-damage claim with 22 weighted checkers across five backends; turns 1--3 are shown here; the remaining three turns follow the same template (wake-up prompt, loud/silent events, per-turn checkers). Each card is one in-universe working day. Coloured pills list the backends the turn touches. Italicised blocks are the turn-entry prompts. Solid dots denote \emph{explicit} events delivered in the prompt; dashed dots denote \emph{silent} mutations injected between turns without a notification (on Wed 5/15, the warehouse sensor log, the overwritten rate table, and the access-log entry all appear without being requested by the agent). Checker rows show per-turn rubric items with weights (\texttt{w1}, \texttt{w1.5}, \texttt{w2}); \textbf{red-line} items (in red) are high-weight hard constraints implemented as deterministic checker failures inside the same weighted rubric (here the day-2 red-line forbids approving or rejecting the claim before the fire-department final report arrives). Scoring is fully rule-based: each checker is a deterministic Python function that inspects post-turn service state.}
  \label{fig:task-card}
\end{figure}
